% MA 214 Discussion 6 -- covers Lecture 8.
%
% Student handout. Page 1 is the concept review; the question set runs from page 2.
% Compile from inside this folder:  pdflatex Discussion6.tex

\documentclass[11pt]{article}
\usepackage[margin=1in]{geometry}
\usepackage{parskip}
\usepackage{fancyhdr}
\usepackage{array}
\usepackage{booktabs}
\usepackage{enumitem}
\usepackage{amsmath}
\usepackage{tabularx}
\usepackage{listings}

\pagestyle{fancy}
\fancyhf{}
\cfoot{\thepage}
\rhead{Discussion 6}
\lhead{MA 214: Applied Statistics}
\renewcommand{\headrulewidth}{.4mm}

\newcommand{\blankline}[1]{\underline{\hspace{#1}}}
\newcolumntype{L}{>{\raggedright\arraybackslash}X}

\lstset{
  basicstyle=\ttfamily\small,
  columns=fullflexible,
  frame=single,
  framerule=0.4pt,
  breaklines=true,
  showstringspaces=false,
  aboveskip=8pt,
  belowskip=8pt
}

% Section header: bold title over a hairline rule.
\newcommand{\parttitle}[1]{%
  \par\medskip\noindent
  {\large\bfseries #1}\par
  \vspace{-.5em}\noindent\rule{\linewidth}{0.4pt}\par\smallskip}

% Writing space.
\newcommand{\answerspace}[1]{\vspace{#1}}

% Flush-left question list: the label runs in at the margin and the body wraps
% back to the margin, so nothing is pushed far to the right.
\newlist{questions}{enumerate}{1}
\setlist[questions]{label=\textbf{Question \arabic*.}, wide=0pt, itemsep=1.4em, topsep=.8em}

\newlist{parts}{enumerate}{1}
\setlist[parts]{label=\textbf{(\Alph*)}, leftmargin=1.6em, labelsep=.45em,
  itemsep=.7em, topsep=.5em}

\begin{document}

\noindent\textbf{Name:}\blankline{2.3in}\hfill\textbf{BUID:}\blankline{1.6in}

\begin{center}
  {\Large\bfseries Discussion 6: What Exactly Gets Shuffled}
\end{center}

\noindent\textbf{Goals:} \textbf{(1)} say what the null distribution is a distribution \emph{of},
\textbf{(2)} match a null hypothesis to the shuffling operation that encodes it, \textbf{(3)} turn
simulation counts into one- and two-sided $p$-values and $s$-values, \textbf{(4)} judge how many
simulations are enough, \textbf{(5)} name the two decision errors in context.

\parttitle{Part 1: Concept Review}

{\small

\noindent Everything in Part 2 can be answered from this page. Keep it in front of you.

\begin{enumerate}[label=\textbf{\arabic*.}, wide=0pt, itemsep=.45em, topsep=.5em]

\item \textbf{Three distributions, three meanings of ``one dot.''} In the distribution of the
\textbf{data}, one dot is one observation. In the \textbf{sampling distribution} of a statistic,
one dot is one whole sample. In the \textbf{null distribution}, one dot is one whole
\emph{simulated} data set, built so that the null hypothesis is true by construction. The first is
observable, the second usually is not, and the third you can always generate.

\item \textbf{The procedure.} State $H_0$ and $H_A$. Choose a test statistic. Build the null
distribution of that statistic. Locate the observed value in it. Report the $p$-value, then decide.

\item \textbf{The $p$-value and the $s$-value.} The $p$-value is the probability of a statistic
\emph{at least as extreme} as the observed one, computed assuming $H_0$. From a simulation with $B$
draws,
\[
p \;\approx\; \frac{\#\{\text{draws at least as extreme}\}}{B},
\qquad
s = -\log_2 p \;\;\text{bits}.
\]
The $s$-value says how many consecutive coin flips coming up heads would be equally surprising:
$p=0.05$ is about $4.3$ bits, $p=0.01$ is about $6.6$ bits.

\item \textbf{What you shuffle \emph{is} the null hypothesis.} The physical operation encodes what
you are willing to assume.

\begin{center}
\renewcommand{\arraystretch}{1.15}
\begin{tabular}{lll}
\toprule
\textbf{Question} & \textbf{$H_0$} & \textbf{What you shuffle or simulate} \\
\midrule
One proportion vs a claim & $\theta = \theta_0$ & nothing; draw from Binomial$(n,\theta_0)$ \\
Two groups, compare means & no group difference & the \emph{group labels} \\
Paired before and after & no within-pair effect & the \emph{sign} of each pair's difference \\
Two categorical variables & independence & one variable's labels \\
\bottomrule
\end{tabular}
\end{center}

\item \textbf{How many simulations.} The simulated $p$ is itself an estimate, with Monte Carlo
standard error
\[
\text{MCSE} = \sqrt{\frac{p(1-p)}{B}} .
\]
Quartering the error takes $16$ times the simulations. A simulated $p$ is never exactly $0$; report
$(\text{count}+1)/(B+1)$, or just say $p < 1/B$.

\item \textbf{Decision errors.} A \textbf{Type I} error rejects a true $H_0$; its long-run rate is
$\alpha$, which you choose. A \textbf{Type II} error fails to reject a false $H_0$. Lowering
$\alpha$ trades Type I risk for Type II risk. Which one to fear depends on what each mistake costs.

\item \textbf{Four traps.} Saying the null distribution is centered at the observed statistic; it
is centered where $H_0$ puts it. Reporting a one-sided $p$ for a two-sided question. Treating
$p = 0.05$ as a cliff. Believing more simulations can change the answer to a question about the
data rather than about the simulation.

\end{enumerate}

}

\newpage

\parttitle{Part 2: Question Set}

\noindent\textbf{Scenario.} The library system tried a Saturday-hours program at $12$ of its
branches, chosen at random; $12$ comparable branches kept their old schedule. The outcome is the
change in weekly visits, in hundreds, over the following term.

\begin{center}
\renewcommand{\arraystretch}{1.15}
\begin{tabular}{ll}
\toprule
\textbf{Program (12 branches)} & 4.1 \; 2.8 \; 3.6 \; 3.9 \; 1.7 \; 4.0 \; 2.2 \; 3.1 \; 2.4 \; 3.4 \; 1.6 \; 4.4 \\
\textbf{Old schedule (12)} & 1.9 \; 3.2 \; 0.8 \; 2.6 \; 1.1 \; 3.7 \; 2.1 \; 0.5 \; 2.9 \; 1.4 \; 3.0 \; 2.2 \\
\bottomrule
\end{tabular}
\end{center}

\noindent Group means are $3.100$ and $2.117$, so the observed difference is
$\bar{x}_{\text{prog}} - \bar{x}_{\text{old}} = 0.983$. The question the director asks is whether
the program \textbf{changed} visits.

\begin{questions}

%%%%%%%%%%%%%%%%%%%%%%%%%%%%%%%%%%%%%%
\item \textbf{What the null distribution is a distribution of.} You shuffle the 24 branch labels
$B = 1000$ times, each time recomputing the difference in group means. The resulting null
distribution has mean $-0.012$ and standard deviation $0.438$.

\begin{parts}
\item What is one dot in that histogram? Be precise: it is not one branch and not one group.
\answerspace{3.5em}

\item The null distribution is centered essentially at $0$. Explain what forces that, and why it
is \emph{not} evidence that the program did nothing.
\answerspace{4em}

\item Suppose the program had been twice as effective, so the observed difference were $1.97$.
Which moves more, the observed statistic or the center of the null distribution? Name the one
feature of the null distribution that stays the same no matter how effective the program was, and
say what guarantees it.
\answerspace{4.5em}

\item How many standard deviations from the center does the observed $0.983$ sit? What does that
number tell you before you compute any $p$-value?
\answerspace{3.5em}
\end{parts}

%%%%%%%%%%%%%%%%%%%%%%%%%%%%%%%%%%%%%%
\item \textbf{Design the shuffle.} Four different questions the library could ask. For each, state
$H_0$ in words or symbols, and say exactly what you would shuffle or simulate. One of them involves
no shuffling at all.

\begin{center}
\renewcommand{\arraystretch}{2.2}
\begin{tabularx}{\linewidth}{LLL}
\toprule
\textbf{Question} & \textbf{$H_0$} & \textbf{What you shuffle or simulate} \\
\midrule
Did the Saturday program change mean visits? & & \\
Is the share of branches meeting target really $0.40$, as the board claims? & & \\
For the same 12 branches, did visits change from fall to spring? & & \\
Does meeting the target depend on branch type? & & \\
\bottomrule
\end{tabularx}
\end{center}

\begin{parts}
\item Two of the four rows shuffle a \emph{label}, and they are not the same label. Say which is
shuffled in each.
\answerspace{4em}

\item In the third row, why is shuffling the group labels the wrong operation?
\answerspace{4em}
\end{parts}

%%%%%%%%%%%%%%%%%%%%%%%%%%%%%%%%%%%%%%
\item \textbf{From counts to a decision.} Of the $1000$ shuffles, \textbf{10} gave a difference
of $0.983$ or larger, and \textbf{24} gave an absolute difference of $0.983$ or larger.

\begin{parts}
\item Give the one-sided $p$-value and the two-sided $p$-value. Which one answers the director's
question as stated, and why?
\answerspace{4em}

\item Compute the $s$-value for the two-sided $p$. Interpret it as a number of coin flips.
\answerspace{3.5em}

\item If the null distribution were perfectly symmetric, the two-sided count would be exactly
twice the one-sided count. It is $24$, not $20$. What does that discrepancy come from, and should
it worry you?
\answerspace{4.5em}

\item Suppose instead that \emph{zero} of the 1000 shuffles had been at least as extreme. Why
would it be wrong to report $p = 0$, and what would you report?
\answerspace{4em}
\end{parts}

%%%%%%%%%%%%%%%%%%%%%%%%%%%%%%%%%%%%%%
\item \textbf{Enough simulations, and the two errors.} Use $p = 0.024$ throughout.

\begin{parts}
\item Compute the Monte Carlo standard error of that $p$ at $B = 1000$. Then give it for
$B = 100$ and $B = 10{,}000$.
\answerspace{4.5em}

\item A classmate reports $p = 0.024$ from $B = 100$ shuffles and concludes the program worked.
Using your answer to (A), explain what is fragile about that.
\answerspace{4em}

\item The board will fund the program city-wide if the test is significant. Describe the Type I
error and the Type II error \emph{in this context}, in terms of money and readers.
\answerspace{4.5em}

\item Which error would you rather make here, and what would you do to $\alpha$ as a result?
\answerspace{3.5em}
\end{parts}

%%%%%%%%%%%%%%%%%%%%%%%%%%%%%%%%%%%%%%
\item \textbf{Find the bugs.} A classmate writes the two-sided permutation test below.
\texttt{all} holds all 24 outcomes, with the 12 program branches first. The code runs, prints a
plausible number, and is wrong in \textbf{three} separate ways.

\begin{lstlisting}[language=R]
obs <- mean(all[1:12]) - mean(all[13:24])

perm <- replicate(20, {
  prog <- sample(all, 12, replace = TRUE)
  old  <- sample(all, 12, replace = TRUE)
  mean(prog) - mean(old)
})

mean(perm > obs)
\end{lstlisting}

\begin{parts}
\item Bug 1 is \texttt{replace = TRUE}. What procedure is this code actually carrying out, and why
does it break the logic of a permutation test? Mention what can happen to a single branch's value.
\answerspace{4.5em}

\item Bug 2 is the number $20$. What is wrong with it, and roughly what is the smallest $p$-value
this code could ever report?
\answerspace{4em}

\item Bug 3 is the last line, given that the director asked whether the program \emph{changed}
visits. Fix it.
\answerspace{3.5em}

\item Rewrite the whole block correctly.
\answerspace{6em}
\end{parts}

\end{questions}

\end{document}
